\chapter*{理论依赖关系}
\addcontentsline{toc}{chapter}{理论依赖关系}

全书按“统计结构 $\to$ 经典连续场 $\to$ 量子动力学”的顺序推进。统计物理部分从有限 Gibbs 系统与配分函数开始，经 Lee--Yang 零点、graph/cycle space、平面对偶、高低温展开、转移算子、Clifford/Jordan--Wigner 与有限扇区到 Onsager--Yang，再由一般 cluster/Mayer/virial 展开和 Beth--Uhlenbeck 回到二体谱与散射数据。

连续介质与经典场部分先建立一般四阶变分算子、边界 Green 配对、自伴实现、零模与 pseudoinverse，再从三维线弹性以 $h/L$ 为控制参数约化到一般 $A$--$B$--$D$ 薄板和 Kirchhoff--Love 特例。随后 Maxwell 方程与守恒四流给出精确辐射谱，世界线、多极与软极限在其后特殊化；当源运动也需自洽求解时，再由 retarded/advanced Green 函数、Dirac 分解和能动量守恒进入 ALD 与受控的 Landau--Lifshitz 降阶。

量子部分先在 $L^2(\mathbb R^d)$ 上建立 $SO(d)$ 分解、径向奇异 Sturm--Liouville 算子、自伴定义域和一般短程散射，再讨论真实 $d$ 维 Maxwell--Poisson 势、固定 $1/r$ 数学模型、零程极限与可解势。最后把封闭系统扩展到系统--环境联合 Hilbert 空间，从精确约化映射依次施加 Born、Markov、coarse-graining/secular 等条件，得到相应适用范围内的 GKSL 动力学。


原子与光场部分不再从二能级或 simple-man 特例反推一般结论。先由 $SU(2)$ 张量积分解和 Wigner--Eckart 定理建立角动量选择结构，并在反对称 $N$ 电子空间中得到 Hartree--Fock 母方程；随后用 Dirac 最小耦合和 Foldy--Wouthuysen 块对角化建立精细结构。一般时间序传播、Dyson--Magnus 与投影记忆核在强场理论之前给出，因此 Fermi 黄金法则、RWA/Rabi、Volkov/SFA 和 Keldysh 复时间鞍点都具有清楚的上游母理论。强场经典回碰先确定 $U_p$ 与 cutoff 尺度，量子强场再解释规范协变、复鞍点与 Floquet 结构；冷却部分则把内部 optical Bloch 稳态与质心 Fokker--Planck 动力学连接起来。

分子与量子光学部分从 Born--Huang 精确通道方程开始，并明确区分完整电子基中的协变动能与有限电子子空间中由正交补空间产生的几何标量势。质量加权 Hessian 给出一般多原子正常模，双原子振转和 Franck--Condon 位移模型随后作为特例；两个二次势能面之间的一般关系由 Duschinsky--Doktorov Gaussian 变换描述。电磁场先作为自伴 Maxwell 正常模系统量子化，再由 Green dyadic 的虚部统一自由空间自发辐射和结构化光学环境。有限群投影算符负责 Hessian 的对称分块，$SO(3)$ 左右作用负责一般分子转子，最后 Schur-complement resolvent 把 Fano--Feshbach 共振与局部 Landau--Zener crossing 放进同一“保留子空间--消去子空间”框架。
